\documentclass{article}
\usepackage{tikz}
\usepackage[paperwidth=15cm, paperheight=16cm, margin=0.5cm]{geometry}
\pagestyle{empty}

\begin{document}

%% Current Question:
\def\Q{1}

%% All 24 Questions:
\expandafter\def\csname Qrow1\endcsname{6/2/3/3}
\expandafter\def\csname Qrow2\endcsname{6/2/2/2}
\expandafter\def\csname Qrow3\endcsname{3/3/2/2}
\expandafter\def\csname Qrow4\endcsname{6/2/4/3}
\expandafter\def\csname Qrow5\endcsname{6/2/7/4}
\expandafter\def\csname Qrow6\endcsname{4/3/7/4}
\expandafter\def\csname Qrow7\endcsname{8/0/4/4}
\expandafter\def\csname Qrow8\endcsname{8/0/2/2}
\expandafter\def\csname Qrow9\endcsname{4/4/2/2}
\expandafter\def\csname Qrow10\endcsname{8/0/4/3}
\expandafter\def\csname Qrow11\endcsname{8/0/1/1}
\expandafter\def\csname Qrow12\endcsname{4/3/1/1}
\expandafter\def\csname Qrow13\endcsname{6/2/3/3}
\expandafter\def\csname Qrow14\endcsname{6/2/2/2}
\expandafter\def\csname Qrow15\endcsname{3/3/2/2}
\expandafter\def\csname Qrow16\endcsname{6/2/4/3}
\expandafter\def\csname Qrow17\endcsname{6/2/7/4}
\expandafter\def\csname Qrow18\endcsname{4/3/7/4}
\expandafter\def\csname Qrow19\endcsname{8/0/4/4}
\expandafter\def\csname Qrow20\endcsname{8/0/2/2}
\expandafter\def\csname Qrow21\endcsname{4/4/2/2}
\expandafter\def\csname Qrow22\endcsname{8/0/4/3}
\expandafter\def\csname Qrow23\endcsname{8/0/1/1}
\expandafter\def\csname Qrow24\endcsname{4/3/1/1}

\expandafter\edef\expandafter\currentrow\expandafter{\csname Qrow\Q\endcsname}

\def\lw{3pt}


%% --- Body colour pool --- UNUSED
%% Update \Nb to (number of colours - 1)
%\def\Nb{0}
%% Human colours
\definecolor{bodycolor0}{HTML}{FFFFFF}
\definecolor{bodycolor1}{HTML}{FFFFFF}
\definecolor{bodycolor2}{HTML}{FFFFFF}
\definecolor{bodycolor3}{HTML}{FFFFFF}

%\definecolor{bodycolor0}{HTML}{FFE8D6}
%\definecolor{bodycolor1}{HTML}{E8A87C}
%\definecolor{bodycolor2}{HTML}{7B3F00}
%\definecolor{bodycolor3}{HTML}{2D1B14}

%% Muppet colours
%\definecolor{bodycolor0}{HTML}{E8233A}
%\definecolor{bodycolor1}{HTML}{1B7FE8}
%\definecolor{bodycolor2}{HTML}{2DB84B}
%\definecolor{bodycolor3}{HTML}{F5A800}


%% --- Shirt colour pool ---
%% Update \Ns to (number of colours - 1)
\def\Ns{23}
\definecolor{shirtcolor0}{HTML}{C8E6A0}
\definecolor{shirtcolor1}{HTML}{A0C8E6}
\definecolor{shirtcolor2}{HTML}{E69830}
\definecolor{shirtcolor3}{HTML}{80FFC8}
\definecolor{shirtcolor4}{HTML}{A0E6D4}
\definecolor{shirtcolor5}{HTML}{A0D4C8}
\definecolor{shirtcolor6}{HTML}{E6E6A0}
\definecolor{shirtcolor7}{HTML}{A0E6C8}
\definecolor{shirtcolor8}{HTML}{E6C8A0}
\definecolor{shirtcolor9}{HTML}{A0A0E6}
\definecolor{shirtcolor10}{HTML}{C8E6E6}
\definecolor{shirtcolor11}{HTML}{C8D4A0}
\definecolor{shirtcolor12}{HTML}{B0D4A0}
\definecolor{shirtcolor13}{HTML}{A0C8B0}
\definecolor{shirtcolor14}{HTML}{B0C8A0}
\definecolor{shirtcolor15}{HTML}{E6D4A0}
\definecolor{shirtcolor16}{HTML}{A0D4E6}
\definecolor{shirtcolor17}{HTML}{D4E6A0}
\definecolor{shirtcolor18}{HTML}{A0C8D4}
\definecolor{shirtcolor19}{HTML}{A0E6A0}
\definecolor{shirtcolor20}{HTML}{B0D4B0}
\definecolor{shirtcolor21}{HTML}{A0B0D4}
\definecolor{shirtcolor22}{HTML}{D4C8A0}
\definecolor{shirtcolor23}{HTML}{A0B8C8}

%% Main drawing command
%% #1=bodyA index, #2=shirtA index, #3=bodyB index, #4=shirtB index
%% #5=tokens for A (0-8), #6=tokens for B (0-8)
\newcommand{\drawImage}[6]{%
\scalebox{1.8}{
  \begin{tikzpicture}[line width=2pt, line cap=round, line join=round]
% No table
%    \draw[fill=brown!25, draw=brown!50]
%      (-6.5, 0.0) -- (6.5, 0.0) -- (6.5, -0.5) -- (-6.5, -0.5) -- cycle;
%    \draw[fill=brown!35, draw=brown!60]
%      (-6.5,-0.5) -- (6.5,-0.5) -- (6.2,-1.1) -- (-6.2,-1.1) -- cycle;
\draw (-3.5, 0.0) -- (3.5, 0.0);
\draw (-3.5, -4.1) -- (3.5, -4.1);

    \player{0.0}{#1}{#2}
 %   \player{ 3.0}{#3}{#4}

    \ifnum#5>0
      \pgfmathsetmacro{\startA}{0 - (#5-1)*0.80/2}
      \foreach \i in {1,...,#5}{
        \pgfmathsetmacro{\tx}{\startA + (\i-1)*0.80}
             \token{\tx}{-0.27}{red}
      }
    \fi
    \ifnum#6>0
      \pgfmathsetmacro{\startB}{0 - (#6-1)*0.80/2}
      \foreach \i in {1,...,#6}{
        \pgfmathsetmacro{\tx}{\startB + (\i-1)*0.80}
             \token{\tx}{-3.77}{red}
      }
    \fi

%    \node[font=\normalsize] at (-3.0, -1.7) {Giocatore A};
%    \node[font=\normalsize] at ( 3.0, -1.7) {Giocatore B};
  \end{tikzpicture}%
}
}


%% How to draw a stick-figure
%\newcommand{\player}[3]{
%  \draw[fill=bodycolor#2, fill opacity=0.7] (#1, 3.6) circle (1.0);
%  \fill (#1-0.32, 3.85) circle (0.14);
%  \fill (#1+0.32, 3.85) circle (0.14);
%  %\draw (#1-0.38, 3.35) arc(220:320:0.50);
%
%  \draw (#1-1.25, 1.65) -- (#1-2, 1.5);
%  \draw (#1+1.25, 1.65) -- (#1+2, 1.5);
%
%  \draw[fill=bodycolor#2, fill opacity=0.7]
%    (#1-0.20, 2.6) -- (#1-0.20, 2.30)
%    .. controls (#1-0.20, 2.10) and (#1+0.20, 2.10) ..
%    (#1+0.20, 2.30) -- (#1+0.20, 2.6);
%
%  \draw[fill=shirtcolor#3, fill opacity=0.7]
%    (#1-0.45, 2.28)
%    -- (#1-1.35, 1.95)
%    -- (#1-1.15, 1.30)
%    -- (#1-0.45, 1.30)
%    -- (#1-0.65, 0.0)
%    -- (#1+0.65, 0.0)
%    -- (#1+0.45, 1.30)
%    -- (#1+1.15, 1.30)
%    -- (#1+1.35, 1.95)
%    -- (#1+0.45, 2.28)
%    .. controls (#1+0.32, 2.36) and (#1+0.22, 2.38) .. (#1+0.20, 2.30)
%    .. controls (#1+0.20, 2.10) and (#1-0.20, 2.10) .. (#1-0.20, 2.30)
%    .. controls (#1-0.22, 2.38) and (#1-0.32, 2.36) ..
%    cycle;
%}
\newcommand{\player}[3]{
  % Trapezoid torso (drawn first so head overlaps it)
  \draw[fill=shirtcolor#3, fill opacity=0.9, line width = 2pt]
    (#1-0.4, 2.2) -- (#1+0.4, 2.2)
    -- (#1+0.75, 0.0) -- (#1-0.75, 0.0) -- cycle;

  % Left arm
  \draw [line width=\lw] (#1-0.48, 1.8)
    .. controls (#1-0.6, 1.6) and (#1-1.2, 1.0) ..
    (#1-1.0, 0.25);
  \draw [line width=\lw] (#1-1.0, 0.25)
	.. controls (#1-0.98,0.1) ..
	(#1 -1.15,0.05);
	
  % Right arm
  \draw [line width=\lw] (#1+0.48, 1.8)
    .. controls (#1+0.6, 1.6) and (#1+1.2, 1.0) ..
    (#1+1.0, 0.25);
  \draw [line width=\lw] (#1+1.0, 0.25)
	.. controls (#1+0.98,0.1) ..
	(#1+1.15,0.05);

  % Head (drawn after torso so it overlaps the top edge)
  \draw[fill=bodycolor#2, fill opacity=1] (#1, 2.9) circle (0.9);
\draw[fill=black] (#1-0.24, 3.15) ellipse (0.08 and 0.14);
\draw[fill=black] (#1+0.24, 3.15) ellipse (0.08 and 0.14);
}

%% Tokens
\newcommand{\token}[3]{% #1=x, #2=y, #3=colour
% Drop shadow
  \draw[fill=black, draw=none, fill opacity = 0.2 ] (#1-0.05, #2-0.07) ellipse (0.38 and 0.18);
    % Shadow
  \draw[fill=#3!60!black, draw=none] (#1, #2-0.04) ellipse (0.38 and 0.18);
  % Main body
  \draw[fill=#3, draw=#3!60!black, line width=0.6pt] (#1, #2) ellipse (0.38 and 0.16);
  % Highlight
  \draw[fill=#3!60!white, draw=none, fill opacity = 0.5] (#1+0.05, #2+0.06) ellipse (0.18 and 0.06);
}

%%\foreach \ta/\tb in {6/2, 3/3, 2/2, 4/3, 7/4, 8/0, 4/4, 2/2, 1/1}{
%\foreach \ta/\tb in {3/3}{
%	\foreach \s in {0,...,\Ns} {
%		\drawImage{0}{\s}{0}{\s}{\ta}{\tb}\clearpage
%	}
%}

\def\runloops#1/#2/#3/#4\relax{%
  \foreach \s in {0,...,\Ns}{
    \drawImage{0}{\s}{0}{\s}{#1}{#2}\clearpage
  }
  \foreach \s in {0,...,\Ns}{
    \drawImage{0}{\s}{0}{\s}{#3}{#4}\clearpage
  }
}
\expandafter\runloops\currentrow\relax

%    \foreach \sa in {0,...,1} {
%      \foreach \ba in {0} {
%        	\foreach \sb in {0,...,1} {
%      	  \foreach \bb in {0,...,1} {
%%    \foreach \sa in {0,...,\Ns} {
%%      \foreach \ba in {0,...,\Nb} {
%%        	\foreach \sb in {0,...,\Ns} {
%%      	  \foreach \bb in {0,...,\Nb} {
%  %        \pgfmathtruncatemacro{\valid}{(\ba < \bb) && (\sa < \sb)}
%        \pgfmathtruncatemacro{\valid}{(\ba != \bb) && (\sa < \sb)}
%        \ifnum\valid=1
%%          \drawImage{\ba}{\sa}{\bb}{\sb}{8}{2}\clearpage
%          \drawImage{\ba}{\sb}{\ba}{\sb}{4}{4}\clearpage
%        \fi
%      }
%    }
%  }
%}

\end{document}